\documentclass[11pt,a4paper]{article}
\usepackage[UTF8]{ctex}
\usepackage{amsmath}
\usepackage{booktabs}
\usepackage{geometry}
\usepackage{xcolor}
\usepackage{hyperref}

\geometry{left=2.5cm,right=2.5cm,top=2.5cm,bottom=2.5cm}

\title{\textbf{Task 2.3: 主题合并与去冗余技术报告}}
\date{2026年2月4日}

\begin{document}

\maketitle

\section{任务概述}

Task 2.3在LDA主题训练之后、跨期对齐之前引入“主题合并”步骤，用于合并高相似度主题、减少冗余与噪音主题，提升后续对齐的稳定性与可解释性。本任务基于主题向量的余弦相似度进行聚类，并输出合并后的主题向量与映射关系。

\section{方法与算法}

\subsection{相似度计算}

对每个时间窗口的主题向量矩阵 $\mathbf{V}\in\mathbb{R}^{K\times d}$，计算主题间余弦相似度：

\begin{equation}
\text{sim}(i,j)=\cos(\mathbf{v}_i,\mathbf{v}_j)=\frac{\mathbf{v}_i\cdot\mathbf{v}_j}{\|\mathbf{v}_i\|_2\,\|\mathbf{v}_j\|_2}
\end{equation}

\subsection{合并策略}

采用基于阈值的贪婪合并策略：
\begin{itemize}
    \item 相似度阈值：0.75
    \item 单簇最大主题数：5
    \item 每窗口最小主题数：20
\end{itemize}

当两个主题相似度高于阈值且合并后簇大小不超过限制时，执行合并。合并后的主题向量为簇内向量均值，并进行L2归一化。

\section{参数配置}

\begin{table}[h]
\centering
\begin{tabular}{ll}
\toprule
\textbf{参数} & \textbf{值} \\
\midrule
相似度阈值 (similarity\_threshold) & 0.75 \\
最小主题数 (min\_topics\_per\_window) & 20 \\
最大簇大小 (max\_cluster\_size) & 5 \\
相似度度量 & 余弦相似度 \\
\bottomrule
\end{tabular}
\caption{主题合并超参数配置}
\end{table}

\section{合并结果汇总}

\subsection{各窗口合并统计}

\begin{table}[h]
\centering
\begin{tabular}{lrrrr}
\toprule
\textbf{时间窗口} & \textbf{原始主题数} & \textbf{合并后主题数} & \textbf{减少率} & \textbf{簇数} \\
\midrule
2016--2017 & 60 & 46 & 23.33\% & 46 \\
2018--2019 & 60 & 50 & 16.67\% & 50 \\
2020--2021 & 60 & 52 & 13.33\% & 52 \\
2022--2023 & 60 & 53 & 11.67\% & 53 \\
2024--2025 & 60 & 55 & 8.33\% & 55 \\
\midrule
\textbf{合计} & 300 & 256 & \textbf{14.67\%} & 256 \\
\bottomrule
\end{tabular}
\caption{主题合并结果统计}
\end{table}

\subsection{关键观察}

\begin{itemize}
    \item 合并力度随时间窗口递增而降低，可能反映后期主题结构更清晰。
    \item 总体主题数从300降至256（减少44个），降低了约14.67\%的冗余。
    \item 合并后主题数仍显著高于最小阈值20，保留足够细粒度。
\end{itemize}

\section{输出文件}

\begin{itemize}
    \item \texttt{output/lda/merged\_topic\_vectors/window\_*\_merged\_vectors.npy}：合并后的主题向量
    \item \texttt{output/lda/merged\_topic\_vectors/window\_*\_cluster\_info.pkl}：簇信息与映射关系
    \item \texttt{output/lda/merged\_topic\_vectors/window\_*\_merged\_labels.pkl}：合并后主题标签
    \item \texttt{output/lda/topic\_merge\_report.csv}：合并统计报告
    \item \texttt{output/lda/topic\_merge\_mapping.pkl}：全局主题映射
\end{itemize}

\section{结论}

Task 2.3通过基于主题向量的相似度合并，有效减少冗余主题并保持语义结构完整，为后续的跨窗口对齐（Task 3.2）提供更稳定的主题集合。总体减少率14.67\%，表明合并策略在降噪与保真之间取得了较好的平衡。

\end{document}
